\documentclass[11pt,a4paper, 
swedish, english %% Make sure to put the main language last!
]{article}
\pdfoutput=1

%% Andréas's custom package 
%% (Will work for most purposes, but is mainly focused on physics.)
\usepackage{../custom_as}

%% Figures can now be put in a folder: 
\graphicspath{ {figures/} %{some_folder_name/}
}

%% If you want to change the margins for just the captions
\usepackage[margin=10 pt]{caption}

%% To add todo-notes in the pdf
\usepackage[%disable  %%this will hide all notes
]{todonotes} 

%% Change the margin in the documents
\usepackage[
%            top    = 3cm,              %% top margin
%            bottom = 3cm,              %% bottom margin
%            left   = 2.5cm, right  = 2.5cm %% left and right margins
]{geometry}


%% If you want to change the formatting of the section headers
%\renewcommand{\thesection}{...}

\newcommand{\EM}{\text{EM}}
\newcommand{\E}{\text{E}}
%\newcommand{\GR}{\text{GR}}
%\newcommand{\SR}{\text{SR}}


%%%%%%%%%%%%%%%%%%%%%%%%%%%%%%%%%%%%%%%%%%%%%%%%%%%%%%%%%%%%%%%%%%%%%%
\begin{document}%% v v v v v v v v v v v v v v v v v v v v v v v v v v
%%%%%%%%%%%%%%%%%%%%%%%%%%%%%%%%%%%%%%%%%%%%%%%%%%%%%%%%%%%%%%%%%%%%%%


%% If you want to use an external file for the title page
%\renewcommand{\thefootnote}{\fnsymbol{footnote}}

%kortkommandon för mailaddresserna
\newcommand{\andsunds}{andsunds@student.chalmers.se}
\newcommand{\rigon}{rigon@student.chalmers.se}



\pagenumbering{roman} %%Romersk sidnumrering i början
\begin{titlepage}
\newgeometry{top=3cm, bottom=2cm}

\newcommand{\HRule}{\rule{\linewidth}{0.5mm}} % Defines a new command for the horizontal lines, change thickness here

\center % Center everything on the page
 
%------------------------------------------------------------------------------------
%	HEADING SECTIONS
%------------------------------------------------------------------------------------

\textsc{\huge Chalmers tekniska högskola}\\[1.5cm] % Name of university/college
\textsc{\Large Rapport, Experimentell fysik 2}\\[0.2cm] % Major heading such as course name
\textsc{\large Termodynamik -- Uppgift 3 }\\[0.5cm] % Minor heading such as course title

%------------------------------------------------------------------------------------
%	TITLE SECTION
%------------------------------------------------------------------------------------

\HRule \\[0.4cm]
{ \LARGE \bfseries 
Studier av kvicksilveratomens atomära emissionsspektra samt absorptionsspektra av laserfärgämnena Rhodamin B och Kumarin 307
}\\[0.4cm] % Title of  document
\HRule \\[1.5cm]
 
%------------------------------------------------------------------------------------
%	AUTHOR SECTION
%------------------------------------------------------------------------------------

\begin{minipage}{0.4\textwidth}
\begin{flushleft} \large
\emph{Författare:}\\
Andréas Sundström\footnotemark{} \\
Rigon Demisai\footnotemark{} 
\end{flushleft}
\end{minipage}
~
\begin{minipage}{0.4\textwidth}
\begin{flushright} \large
\emph{Labassistent:} \\
Martin Wersäll
\end{flushright}
\end{minipage}\\[3cm]

\setcounter{footnote}{0}
\stepcounter{footnote}
  \footnotetext{\href{mailto:\andsunds}{\texttt{\andsunds}}}
\stepcounter{footnote}
  \footnotetext{\href{mailto:\rigon}{\texttt{\rigon}}}



%------------------------------------------------------------------------------------
%	DATE SECTION
%------------------------------------------------------------------------------------
% Följer ISO-standarden för tidsintervall:
% https://en.wikipedia.org/wiki/ISO_8601#Time_intervals
% "Double hyphen" också ok istället för '/'. -- i LaTeX är dock lite på gränsen
{ \large
\begin{tabular}{rc}
    Laboration utförd: & 2015-12-11/15 \\[0.1cm]
    Rapport inlämnad: & \today
\end{tabular}\\[1cm]
}

%------------------------------------------------------------------------------------
%	LOGO SECTION
%------------------------------------------------------------------------------------

\includegraphics[height=5cm]{logo.pdf} % Include a department/university logo
 
%------------------------------------------------------------------------------------

\vfill % Fill the rest of the page with whitespace

\end{titlepage}
\restoregeometry


\setcounter{page}{2}%detta är ANDRA (2) sidan

\renewcommand{\abstractname}{Sammandrag}
\begin{abstract}
Den här rapporten beskriver en spektroskopisk studie av
kvicksilveratomens atomära emissionsspektrum samt en studie av
laserfärgämnena rhodamin~B och kumarin~307 och deras
absorptionsspektra.
Ur emisionspektrumet från Hg har vissa av atomens energinivåer kartlagts,
baserat på kvanmekaniska regler och jämförelse med tidigare data på
Hg.  Kvicksilvrets emissionsspektra
är taget i intervallet 350--1100\,nm, där detekterades totalt 23
emissionstoppar, varefter 13 övergångar kunde identifieras.
Absorptionsspektrumen för rhodamin~B och kumarin~307 visar breda
absorptionsband vilket är kännetecknande för flourescerande ämnen som
består av stora organiska molekyler. 
Mätningarna har utförts med en Spex~270M spektrometer och
datainsamlingen har gjorts i LabView.
\end{abstract}

\renewcommand{\abstractname}{Abstract}
\begin{abstract}
This report describes a spectroscopic study of the atomic emision
spectrum of mercury, and also a study of the laser dyes Rhodamine~D
and Coumarin~307 and their absorption spectrum.
Som of the energylevels of Hg have been idetified from the emission
spectrum, based on quantum mecanichal rules and comparison with erlier
data on Hg. The emission spectrum of mercury was taken in the interval
of 350--1100\,nm, there a total of 23 emission lines were detected,
from which 13 different atomic transitions could be identified. 
The absorption spectrum of Rhodamin~B and Coumarin~307  show wide
absorption bands which are characteristic of big organic molecules.
THe measurements were made with a Spex~270M spectrometer and the data
collection was done through LabVIEW.
\end{abstract}

\clearpage
\renewcommand{\contentsname}{Innehållsförteckning}
\tableofcontents

\clearpage
\pagenumbering{arabic}
\setcounter{page}{1}

\renewcommand{\thefootnote}{\arabic{footnote}}
\setcounter{footnote}{0}



%%%%%%%%%%%%%%%%%%%% vvv Internal title page vvv %%%%%%%%%%%%%%%%%%%%%
\title{Gravitation and Cosmology -- FFM 071
\\ {\Large Hand-in set 1} }
\author{Andréas Sundström}
\date{\today}

\maketitle

%%%%%%%%%%%%%%%%%%%% ^^^ Internal title page ^^^ %%%%%%%%%%%%%%%%%%%%%
%% If you want a list of all todos
%\todolist


\renewcommand{\thesubsection}{\arabic{section} (\alph{subsection})}
\section{Some Lorentz transformations}
\subsection{Proper time}
To show that the proper time is invariant, we just use the definition
of $\rd\tau$ and the transformed coordinates 
$\rd{x'}^\alpha={\Lambda^\alpha}_\gamma \rd{x}^{\gamma}$,
\begin{equation}
\rd{\tau'}^2=-\eta_{\alpha\beta}\rd{x'}^{\alpha}\rd{x'}^{\beta}
=-\eta_{\alpha\beta}({\Lambda^\alpha}_\gamma \rd{x}^{\gamma})
({\Lambda^\beta}_\delta\rd{x}^{\delta})
=-\eta_{\alpha\beta}{\Lambda^\alpha}_\gamma{\Lambda^\beta}_\delta
\rd{x}^{\gamma}\rd{x}^{\delta}.
\end{equation}
next we use the condition that the Lorentz transform must satisfy
\begin{equation}
\eta_{\alpha\beta}{\Lambda^\alpha}_\gamma{\Lambda^\beta}_\delta
=\eta_{\gamma\delta},
\end{equation}
which therefore yields
\begin{equation}
\rd{\tau'}^2=\eta_{\gamma\delta}\rd{x}^{\gamma}\rd{x}^{\delta}
=\rd\tau^2.
\end{equation}
Since we have now concluded that $\rd\tau^2$ is a scalar (invariant)
under Lorentz transforms, so must $\rd\tau=\sqrt{\rd\tau^2}$.

\subsection{Relativistic particle velocity}
For a Lorentz boost with $\vb*v=(v,0,0)$, the transformations are as
follows
\begin{equation}
\begin{cases}
\rd{t}'=\gamma(\rd{t}-v\rd{x})\\
\rd{x}'=\gamma(\rd{x}-v\rd{t})\\
\rd{y}'=\rd{y}\\
\rd{z}'=\rd{z}
\end{cases}
\qand
\begin{cases}
\rd{t}=\gamma(\rd{t}'+v\rd{x}')\\
\rd{x}=\gamma(\rd{x}'+v\rd{t}')\\
\rd{y}=\rd{y}'\\
\rd{z}=\rd{z}'.
\end{cases}
\end{equation}
The velocity of a particle 
\begin{equation}
\vb*u=\qty(\dv{x}{t},\dv{y}{t},\dv{z}{t})
=(u_x,u_y,u_z)=(u\cos\theta,u\sin\theta,0)
\end{equation}
in the first frame, will transform as
\begin{equation}\label{eq1:upx1}
u'_x=\dv{x'}{t'}=\dv{x'}{t}\dv{t}{t'}
=\qty(\pdv{x'}{t}+\pdv{x'}{x}\pdv{x}{t})
\qty(\pdv{t}{t'}+\pdv{t}{x'}\pdv{x'}{t'})
=\gamma^2(-v+u_x)(1+vu'_x)
\end{equation}
and
\begin{equation}\label{eq1:upy1}
u'_y=\dv{y'}{t'}=\dv{y'}{t}\dv{t}{t'}
=\qty(\pdv{y'}{y}\pdv{y}{t})
\qty(\pdv{t}{t'}+\pdv{t}{x'}\pdv{x'}{t'})
=\gamma u_y(1+vu'_x).
\end{equation}
Solving for $u'_x$ in \eqref{eq1:upx1} yields
\begin{equation}
u'_x=\frac{u_x-v}{\gamma^{-2}-v(u_x-v)}
=\frac{u_x-v}{1-v^2-v(u_x-v)}
=\frac{u_x-v}{1-u_x v},
\end{equation}
similarly with \eqref{eq1:upy1}
\begin{equation}
u'_y=\gamma u_y\qty(1+v\frac{u_x-v}{1-u_x v})
=\frac{u_y}{\gamma(1-u_x v)}.
\end{equation}
The transform of $\theta$, the angle $\vb*u$ makes with the $x$ axis,
is easily given by 
\begin{equation}
\tan\theta'=\frac{u_y'}{u_x'}
=\qty(\frac{u\sin\theta}{\gamma(1-u\cos\theta v)})
\qty(\frac{1-u\cos\theta v}{u\cos\theta-v})
=\frac{\sin\theta}{\gamma(\cos\theta-\nicefrac{v}{u})}.
\end{equation}



\section{The stress tensor}
To show the conservation laws for the stress tensor\footnotemark{}
\footnotetext{I call the stress tensor of a particle
  ${T_n}^{\alpha\beta}$, with subscript $n$, and use $x_n(\tau)$ and
  $p_n$ for the trajectory and momentum of the particle, to more
  clearly distinguish what ``belongs'' to the particle and what are
  free variables.} 
\begin{equation}
{T_n}^{\alpha\beta}(x)=\int\rd\tau\,{p_n}^\alpha
\dv{{x_n}^\beta}{\tau}\delta^4\big(x-x_n(\tau)\big).
\end{equation}
The divergence of the stress tensor therefore becomes
\begin{equation}\label{eq2:pdvT1}
\pd_\beta{T_n}^{\alpha\beta}(x)=\int\rd\tau\,{p_n}^\alpha
\dv{{x_n}^\beta}{\tau}\,\pdv{x^\beta}\delta^4\big(x-x_n(\tau)\big).
\end{equation}
We note that the derivatives only act on the delta function since
$x_n$ and $p_n$ are the trajectory and momentum of the particle, while
the derivatives act on the free variable $x$. Next we note that for
any function $f$ 
\begin{equation}
\pdv{x}\Big[f(x-y)\Big]=-\pdv{y}\Big[f(x-y)\Big],
\end{equation}
which means that \eqref{eq2:pdvT1} can be written as
\vspace{-1em}
\begin{equation}
\begin{aligned}
\pd_\beta{T_n}^{\alpha\beta}(x)=&-\int\rd\tau\,{p_n}^\alpha
\overbrace{\dv{{x_n}^\beta}{\tau}\,\pdv{{x_n}^\beta}}^{=\dv*\tau}
\delta^4\big(x-x_n(\tau)\big)\\
=&-\int\rd\tau\,{p_n}^\alpha\dv{\tau}\delta^4\big(x-x_n(\tau)\big).
\end{aligned}
\end{equation}
Integration by parts now yields
\begin{equation}
\begin{aligned}
\pd_\beta{T_n}^{\alpha\beta}(x)=&
-\int\rd\tau\,{p_n}^\alpha\dv{\tau}\delta^4\big(x-x_n(\tau)\big)\\
=&-\cancel{\qty[{p_n}^\alpha\delta^4\big(x-x_n(\tau)\big)]}
+\int\rd\tau\,\dv{{p_n}^\alpha}{\tau}\delta^4\big(x-x_n(\tau)\big).
\end{aligned}
\end{equation}
The first term from the integration by parts cancel since the
integration is over ``all $\tau$'' and therefore $x_n\neq x$ at the
end points, which means that the delta function is zero.

Now if the particle interacts with an electromagnetic field, it will
experience a Lorentz force
\begin{equation}
\dv{{p_n}^\alpha}{\tau}=e_n\eta_{\beta\gamma}F^{\alpha\beta}
\dv{{x_n}^\gamma}{\tau}={F^{\alpha}}_\gamma \,e_n{u_n}^\gamma.
%={F^{\alpha}}_\gamma J^\gamma.
\end{equation}
This means that
\begin{equation}\label{eq2:pdvT}
\pd_\beta{T_n}^{\alpha\beta}(x)=\int\rd\tau\,{F^{\alpha}}_\gamma e_n{u_n}^\gamma
\delta^4\big(x-x_n(\tau)\big)
={F^{\alpha}}_\gamma e_n{u_n}^\gamma
\delta^3\big(\vb*x-\vb*x_n(\tau)\big)
={F^{\alpha}}_\gamma J^\gamma
\end{equation}
which in general is non-vanishing. Note that since we are assuming
that this particle is the only particle the current is indeed
$J^\alpha=e_n{u_n}^\alpha\delta^3\big(\vb*x-\vb*x_n(\tau)\big)$, and
if we were to have more particles all the above calculations would be
exactly the same (but with summation signs in front of them). 

We now take a look at the electromagnetic stress tensor
\begin{equation}
{T_{\EM}}^{\alpha\beta}={F^{\alpha}}_\gamma F^{\beta\gamma}
-\frac{1}{4}\eta^{\alpha\beta}F_{\gamma\delta}F^{\gamma\delta},
\end{equation}
and then Maxwell's equations will come in handy later
\begin{subequations}
\begin{equation}\label{eq:Maxwell}
\pd_\alpha F^{\alpha\beta}=-J^\beta,
\end{equation}
\begin{equation}\label{eq:Bianchi}
%\pd_{[\alpha}F_{\beta\gamma]}=
\pd_\alpha F_{\beta\gamma} +\pd_\beta F_{\gamma\alpha} +\pd_\gamma F_{\beta\alpha} 
=0.
\end{equation}
\end{subequations}
Its divergence therefore is
\begin{equation}
\begin{aligned}
\pd_{\beta}{T_{\EM}}^{\alpha\beta}=&
{F^{\alpha}}_\gamma \qty(\pd_\beta F^{\beta\gamma})
+\qty(\pd_\beta{F^{\alpha}}_\gamma) F^{\beta\gamma}
-\frac{1}{4}\eta^{\alpha\beta}\pd_\beta\qty[F_{\gamma\delta}F^{\gamma\delta}]\\
=&-{F^{\alpha}}_\gamma J^\gamma
+F^{\beta\gamma} \,\pd_\beta{F^{\alpha}}_\gamma
-\frac{1}{2}F_{\gamma\delta}\,\pd^\alpha F^{\gamma\delta}.
\end{aligned}
\end{equation}
The first term is given directly from \eqref{eq:Maxwell}.
We now take a look at the last two terms. Firstly we can raise and
lower the betas and gammas which are contracted in the second term, 
$F^{\beta\gamma} \,\pd_\beta{F^{\alpha}}_\gamma
=F_{\beta\gamma} \,\pd^\beta F^{\alpha\gamma}$, then in the last term
we are free to rename the two contracted indices $\gamma\to\beta$ and
$\delta\to\gamma$. We now have
\begin{equation}
F_{\beta\gamma} \,\pd^\beta F^{\alpha\gamma}
-\frac{1}{2}F_{\beta\gamma}\,\pd^\alpha F^{\beta\gamma}=
\frac{1}{2}\qty( F_{\beta\gamma}\pd^\beta F^{\alpha\gamma}
-F_{\beta\gamma}\pd^\beta F^{\gamma\alpha})
-\frac{1}{2}F_{\beta\gamma}\pd^\alpha F^{\beta\gamma}.
\end{equation}
Here we used the anti-symmetry of the EM tensor, 
$F^{\alpha\gamma}=-F^{\gamma\alpha}$. Now we do some more index
gymnastics by interchanging $\beta\leftrightarrow\gamma$ in the first
term, $F_{\gamma\beta}\pd^\gamma F^{\alpha\beta}
=-F_{\beta\gamma}\pd^\gamma F^{\alpha\beta}$. 
In the end what we have is
\begin{equation}
F_{\beta\gamma} \,\pd^\beta F^{\alpha\gamma}
-\frac{1}{2}F_{\beta\gamma}\,\pd^\alpha F^{\beta\gamma}=
-\frac{1}{2}F_{\beta\gamma}\qty( 
\pd^\gamma F^{\alpha\beta}
+\pd^\beta F^{\gamma\alpha}
+\pd^\alpha F^{\beta\gamma}).
\end{equation}
The expression inside the parentheses is just the contravariant version
of the RHS of \eqref{eq:Bianchi} and therefore vanishes.
So what we are left with for the divergence is
\begin{equation}
\pd_{\beta}{T_{\EM}}^{\alpha\beta}=-{F^{\alpha}}_\gamma J^\gamma,
\end{equation}
which is exactly $-\pd_\beta{T_n}^{\alpha\beta}$ from \eqref{eq2:pdvT}.

Therefore since $\pd_\beta$ is a linear operator we know that
\begin{equation}
\pd_{\beta}\qty[{T_{\EM}}^{\alpha\beta}+{T_n}^{\alpha\beta}]=0.
\end{equation}
This is physically reasonable since the total stress tensor, which
represents the total energy and momentum-current density, should
reasonably be conserved in a system. The two terms
${T_{\EM}}^{\alpha\beta}$ and ${T_n}^{\alpha\beta}$ are not conserved
separately since the particle and EM field exchanges energy and
momentum with each other, which also leads to the above conclusion
that there is no \emph{net} change in the total energy or momentum of
the system.





\renewcommand{\thesubsection}{\arabic{section} (\roman{subsection})}
\section{Time dilation}
\swapcommands{\phi}{\varphi}
\subsection{Two separate effects}
For a satellite in circular orbit around the Earth, 
%as in \figref{fig:orbit}, 
the time dilation can be seen as two separate
effects. Firstly the slowing down of time from the orbital speed due
to special relativity, and secondly a speed-up from weaker gravity due
to general relativity. It is not really correct to take these two as
separate effects, but in a case like the earth where both the orbital
speed and gravitation is extremely small (in natural units) this
approximation holds.

We can begin with the effects of gravity. The strength of gravity in
the weak field limit is given by the Newtonian potential
\begin{equation}\label{eq3:phi}
\phi(r)=-\frac{\mu}{r},
\end{equation}
where $\mu=GM$, $G$ being the constant of gravity and $M$ the mass
of the Earth. For a distant observer, unaffected by Earth's gravity,
a stationary clock (both the Earth and the clock being stationary in
the observer's frame of reference) only the 00 component of the metric
affects the time dilation, since the gravitational field is static in
this frame of reference, this component is in the weak field limit
$g_{00}=-(1+2\phi)$. The time dilation therefore becomes 
\begin{equation}
\rd\tau=
\overbrace{(-g_{\mu\nu}\rd{x}^\mu\rd{x}^\nu)^{1/2}}
^{=(-g_{00})^{1/2}\rd{t},\text{ static clock and field}}
=\Big(1+2\phi(r)\Big)^{1/2}\rd{t}
\approx\qty(1-\frac{\mu}{r})\rd{t},
\end{equation}
where $\rd{x}^\mu$ and $\rd{t}$ are in the coordinate system of the distant
observer, and $\rd\tau$ is the proper time recorded by that stationary
clock. Note that since $\rd\tau<\rd{t}$, the distant observer
will see the clock tick \emph{slower}\footnotemark{} by the amount
given above.  
\footnotetext{Say the clock is ticking with a time $\rd\tau$ between
  ticks. The distant observer sees the clock above Earth tick with a
  time $\rd{t}$ between ticks in his or her frame of reference.
  Therefore since $\rd{t}>\rd\tau$ the Earth clock seems slower than
  an identical clock at the distant observer.}

%\begin{figure}\centering
%\resizebox{.8\textwidth}{!}{
%\input{figures/Laika.pdf_t}%}
%\caption{A circular orbit around the Earth}
%\label{fig:orbit}
%\end{figure}

To get the SR time dilation, we need the orbital speed which is
handily calculated using Newtonian mechanics to be
\begin{equation}
v^2=\frac{GM}{R}=\frac{\mu}{R},
\end{equation}
where $R$ is the orbital radius. Only accounting for the SR effects of
the orbital speed that time dilation becomes
\begin{equation}
\rd\tau=\gamma^{-1}\,\rd{t}=\qty(1-v^2)^{1/2}\rd{t}
=\qty(1-\frac{\mu}{R})^{1/2}\rd{t}
\approx\qty(1-\frac{\mu}{2R})\rd{t}.
\end{equation}
Once again $\rd\tau$ is the proper time of a clock moving at a speed
$v$, without gravity, while $\rd{t}$ is its time recorded by the
distant observer. 

In this weak field and slow moving limit the real proper time of the
satellite should therefore be the product of the two effects
\begin{equation}
\rd\tau_\text{sat}=\qty(1-\frac{\mu}{R})\qty(1-\frac{\mu}{2R})\rd{t}
\approx\qty(1-\frac{3\mu}{2R})\rd{t}.
\end{equation}
This we want to compare to a stationary clock on the surface of the
Earth, at a distance $R_\E$ from the center of the Earth, which only
has the GR effects 
\begin{equation}
\rd\tau_\E=\eval{\rd\tau}_{R_\E}=\qty(1-\frac{\mu}{R_\E})\rd{t}.
\end{equation}
Setting these two equal, we get
\begin{equation}\label{eq3:R}
R=\frac{3}{2}R_\E.
\end{equation}


\subsection{But there is really only one time dilation}
Since the previous section felt a bit hand-wavy, we can also redo the
same calculations a bit more rigorously. We actually had the answer in
front of our noses from the very beginning. We start off with the
square of the proper time compared to the time of the distant observer 
\begin{equation}\label{eq3:dtaudt}
\qty(\frac{\rd\tau}{\rd{t}})^2=
-g_{\mu\nu}\frac{\rd{x}^\mu}{\rd{t}}\frac{\rd{x}^\nu}{\rd{t}}.
\end{equation}
Next, in the weak field approximation we can write the metric as
\begin{equation}
g_{\mu\nu}=\eta_{\mu\nu}+h_{\mu\nu}\qc h_{\mu\nu}\ll1.
\end{equation}
Since we know that the satellite is slow,
\begin{equation}
1=\dv{x^0}{t}\gg\dv{x^i}{t}= v^i
\end{equation}
we know that the only correction to the metric which survives to first
order is $h_{00}=-2\phi$, the rest of them get multiplied by a factor
$v^i\sim v\ll1$.

We can now evaluate \eqref{eq3:dtaudt} to be
\begin{equation}
\qty(\frac{\rd\tau}{\rd{t}})^2\approx
-g_{00}+\eta_{ij}\frac{\rd{x}^i}{\rd{t}}\frac{\rd{x}^j}{\rd{t}}
=1+2\phi-v^2.
\end{equation}
For the orbiting satellite this becomes
\begin{equation}
\qty(\frac{\rd\tau_\text{sat}}{\rd{t}})^2
=1-2\frac{\mu}{R}-\frac{\mu}{R}
=1-3\frac{\mu}{R},
\end{equation}
and for the stationary clock on the surface of the Earth we get
\begin{equation}
\qty(\frac{\rd\tau_\text{E}}{\rd{t}})^2=1-2\frac{\mu}{R}.
\end{equation}
For these two proper times to be the same, we recover \eqref{eq3:R}
from the previous calculation.


\section{Stereographic projection}
\swapcommands{\phi}{\varphi}

\begin{figure}\centering
%\resizebox{.8\textwidth}{!}{
\input{figures/stereographic.pdf_t}%}
\caption{}
\label{fig:stereographic}
\end{figure}

\figref{fig:stereographic} shows the geometry of a stereographic
projection from the south pole of the unit sphere onto the equatorial
plane, for two points on the sphere with the same polar angle
$\phi$. From the Theorem of inscribed angles it is clear that the
vertical angle from the south pole is exactly half the azimuthal
angle $\theta$. This means that the distance from $(x,y)$ to the
origin, $r$, must be given by $r=\tan(\theta/2)$. Therefore,
expressing $x$ and $y$ in polar coordinates we get
\begin{equation}
\begin{cases}
x=\tan\tfrac{\theta}{2}\,\cos\phi\\
y=\tan\tfrac{\theta}{2}\,\sin\phi.
\end{cases}
\end{equation}

Here we want to find the metric 
\begin{equation}
\rd{s}^2=\rd\theta^2+\sin^2\theta\,\rd\phi^2
\end{equation}
in terms of the $(x,y)$ coordinates. As a good first start, we can
begin to compute
\begin{equation}
\rd{x}=\pdv{x}{\theta}\rd\theta+\pdv{x}{\phi}\rd\phi
=\frac{\cos\phi\;\rd\theta}{2\cos^2(\theta/2)}
-\tan\frac{\theta}{2}\,\sin\phi\,\rd\phi
\end{equation}
and
\begin{equation}
\rd{y}=\pdv{y}{\theta}\rd\theta+\pdv{y}{\phi}\rd\phi
=\frac{\sin\phi\;\rd\theta}{2\cos^2(\theta/2)}
+\tan\frac{\theta}{2}\,\cos\phi\,\rd\phi.
\end{equation}
From here we easily calculate
\begin{equation}
\rd{x}^2+\rd{y}^2=
\frac{\rd\theta^2}{4\cos^4(\theta/2)} 
+\tan^2\frac{\theta}{2}\;\rd\phi^2
\end{equation}
by using $\cos^2+\sin^2\equiv1$ and seeing that the cross
terms cancel. To simplify this, we use the trigonometric identities
\begin{equation}
\cos^2(\theta/2)\equiv\frac{1+\cos\theta}{2}
\end{equation}
and
\begin{equation}
\tan\frac{\theta}{2}\equiv \frac{\sin\theta}{1+\cos\theta},
\end{equation}
which gives
\begin{equation}
\rd{x}^2+\rd{y}^2
=\frac{\rd\theta^2+\sin^2\theta\,\rd\phi^2}{(1+\cos\theta)^2} 
=\frac{\rd{s}^2}{(1+\cos\theta)^2}.
\end{equation}
We are almost there. We just need to use one more trigonometric
identity,
\begin{equation}
r^2=\tan^2\frac{\theta}{2}\equiv
\frac{1-\cos\theta}{1+\cos\theta}
=\frac{1+1-(1+\cos\theta)}{1+\cos\theta}
=\frac{2}{1+\cos\theta}-1.
\end{equation}
Solving for $(1+\cos\theta)$ yields
\begin{equation}
1+\cos\theta=\frac{2}{1+r^2},
\end{equation}
and we see that we can write down the metric for this coordinate
system as
\begin{equation}
\rd{s}^2=4\,\frac{\,\rd{x}^2+\rd{y}^2\,}{(1+r^2)^2},
\end{equation}
with $r^2=x^2+y^2$.

\todo[inline]{There is a part ii and iii to this problem as well.}


\section{Variational method}

\subsection{The basic Lagrangian}

\subsection{Another Lagrangian}
\begin{equation}\label{eq5:Lagrangian2}
S'[x]=\int_A^B\rd\tau\,\qty(-g_{\mu\nu}\dv{x^\mu}{\tau}\dv{x^\nu}{\tau})
\end{equation}



\section{Covariant Laplacian}
\swapcommands{\phi}{\varphi}
The covariant Laplacian acting on a scalar field is given by
\begin{equation}
\Box\phi=\nabla_\mu\nabla^\mu\phi=\nabla_\mu(g^{\mu\nu}\nabla_\nu)\phi
=(g^{\mu\nu}\nabla_\mu)\nabla_\nu\phi
=\nabla^\nu\nabla_\nu\phi.
\end{equation}
We are free to move around the metric tensor as we wish because of the
\emph{metric postulate}.

Let's use the middle expression to write out the Laplacian
explicitly. First of all, we note that $\nabla_\nu\phi=\pd_\nu\phi$
since there is no indices for the affine connection to work on. Next
we get
\begin{equation}
\Box\phi=\pd_\mu(g^{\mu\nu}\pd_\nu)\phi
+\Gamma^\mu_{\mu\rho}(g^{\rho\nu}\pd_{\nu}\phi).
\end{equation}
From here we use the determinant formula
\begin{equation}
\Gamma^\mu_{\mu\rho}=\frac{1}{2}
\Big(\cancel{g^{\mu\sigma}\pd_\mu g_{\rho\sigma}}
-\cancel{g^{\mu\sigma}\pd_\sigma g_{\mu\rho}} \Big)
+\frac{1}{2}g^{\mu\sigma}\pd\rho g_{\mu\sigma}
=\frac{1}{2}g^{-1}\pd_\rho g,
\end{equation}
where $g=\det(g_{\mu\nu})$ is the determinant of the metric.

In the even more concrete example of polar coordinates on a flat 2D
surface, we have the metric
\begin{equation}
g_{\mu\nu}=
\begin{bmatrix}
1&0\\
0&r^2
\end{bmatrix}
\qand
g^{\mu\nu}=
\begin{bmatrix}
1&0\\
0&r^{-2}
\end{bmatrix},
\end{equation}
resulting in $g=r^2$ and
\begin{equation}
\Gamma^\mu_{\mu\rho}=\frac{1}{2}g^{-1}
\begin{bmatrix}
\pd_r g\\\pd_\varphi g
\end{bmatrix}
=
\begin{bmatrix}
r^{-1}\\0
\end{bmatrix}.
\end{equation}
We can now write down the covariant Laplacian on a flat surface with
polar coordinates as
\begin{equation}
\Box\phi=\pd_r(g^{rr}\pd_r\phi)
+\pd_\varphi(g^{\varphi\varphi}\pd_\varphi\phi)
+\Gamma^{\mu}_{\mu r}(g^{rr}\pd_r\phi)
=\Big[\pd_{rr}+r^{-1}\pd_r+r^{-2}\pd_{\varphi\varphi}
\Big]\phi,
\end{equation}
which is precisely what we expected from previous knowledge on
multi-variable calculus. 


\section{A geodesic}
To find the geodesics of the two dimensional metric
\begin{equation}
\rd\tau^2=t^{-2}\,\rd{t}^2 - t^{-2}\,\rd{x}^2,
\end{equation}
we begin by noting that it only depends on $t$. This means that if we
use the Lagrangian from \eqref{eq5:Lagrangian2}
\begin{equation}
\mathcal{L}=-g_{\mu\nu}\dv{x^\mu}{\tau}\dv{x^\nu}{\tau}
=t^{-2}\,\dot{t}^2 - t^{-2}\,\dot{x}^2,
\end{equation}
then we can use the Euler-Lagrange equation
\begin{equation}
\dv{\tau}\qty[\frac{\delta\mathcal{L}}{\delta\dot{x}}]
-\cancelto{0}{\frac{\delta\mathcal{L}}{\delta{x}}}=0.
\end{equation}
This means that
\begin{equation}
\dv{\tau}\qty[-2t^{-2}\dot{x}]=0
\end{equation}
or in other words
\begin{equation}
\dot{x}=t^2C,
\end{equation}
where $C$ is a constant. We can also write
\begin{equation}
\dot{t}=\dv{t}{\tau}=\dv{t}{x}\dv{x}{\tau}
=\dv{t}{x}\,t^2C.
\end{equation}
Now, for a time-like ($\rd{t}>\rd{x}$) geodesic we get
\begin{equation}
+1=t^{-2}\,\qty(\dv{t}{\tau})^2
- t^{-2}\,\qty(\dv{x}{\tau})^2
=t^{-2}\,\qty(t^2C\dv{t}{x})^2
- t^{-2}\qty(t^2C)^2
\end{equation}
which leads to
\begin{equation}
\qty(\dv{t}{x})^2=\frac{1+C^2t^2}{C^2t^2}
\quad\Longrightarrow\quad
\dv{x}{t}=\frac{Ct}{\sqrt{1+C^2t^2}}<1\qfor t>0.
\end{equation}
This has the solution
\begin{equation}
x(t)=\sqrt{1+C^2t^2}+x_0,
\end{equation}
for some constant of integration $x_0$. If we consider a light cone
emanating from the origin $(t,x)=(0,0)$, and since we want the
geodesic to lie inside the light cone, $x_0=-1$ and we see that indeed
$x(t)<t$ for $t>0$.








\section{Some curved geometries}








%%%%%%%%%%%%%%%%%%%%%%%%%%%%%%%%%%%%%%%%%%%%%%%%%%%%%%%%%%%%%%%%%%%%%%
\end{document}%% ^ ^ ^ ^ ^ ^ ^ ^ ^ ^ ^ ^ ^ ^ ^ ^ ^ ^ ^ ^ ^ ^ ^ ^ ^ ^ ^
%%%%%%%%%%%%%%%%%%%%%%%%%%%%%%%%%%%%%%%%%%%%%%%%%%%%%%%%%%%%%%%%%%%%%%



%%% Local Variables:
%%% mode: latex
%%% TeX-master: t
%%% End:

%  LocalWords:  contravariant Stereographic covariant variational
%  LocalWords:  azimuthal affine
